\section{Threats to Validity}
\label{sec:threats}

\textbf{FPR--TPR tradeoff.}
B2 reduces CV(FPR) but increases CV(TPR) ($0.336$ vs.\ $0.289$) and decreases P10 Macro-F1
($0.300$ vs.\ $0.344$); worst-client balanced accuracy is not materially degraded.
In this dataset, the P10 Macro-F1 degradation is concentrated in the low-separability
Ennio Doorbell client (benign-vs-attack AUROC $0.931$, lowest among all nine clients),
so the result is best interpreted as a client-specific failure mode of percentile
calibration rather than evidence of a universal FPR--TPR tradeoff.
Per-client threshold personalization improves false-alarm dispersion fairness, not universal detection quality.

\textbf{Dataset scope.}
Regime~A uses 9 physical IoT devices with natural device-level heterogeneity
(confirmatory); Regime~C ($K{=}20$ synthetic clients) tests severity variation;
Regime~B provides a contrasting near-homogeneous applicability boundary.
The results are a controlled measurement of threshold-scope effects under the tested
protocol, not fleet-scale deployment validation; larger multi-vendor fleets are future work.

\textbf{CICIoT2023 partition.}
Regime~B uses file-level pseudo-clients with randomized within-file benign ordering;
no temporal drift is preserved.
The null result is therefore an applicability boundary under this randomized near-homogeneous
partition only---not evidence that threshold personalization is unnecessary for CICIoT2023
generally.
Chronological, physical-device, attack-source, or attack-type CICIoT2023 partitions would define a different heterogeneity regime and remain outside the present threshold-only controlled comparison.

\textbf{Calibration data sufficiency.}
In all reported conditions, calibration-pending count is zero.
In deployments with fewer than $n_{\min}=100$ benign calibration samples per client,
B2 and B4 fall back to $\tau_{\text{global}}$, reducing their advantage.

\textbf{Threshold quantile sensitivity.}
$q$ is fixed at $0.95$. A sensitivity check over $q \in \{0.90, 0.95, 0.99\}$
preserves the qualitative B2 $<$ B4 $<$ B1 ordering.
The B1-pooled and B1-weighted sensitivity checks are consistent with
the B1-to-B2 gap not being driven solely by arithmetic-mean construction (Table~\ref{tab:sensitivity_shared}).

\textbf{Attack scope and adversarial robustness.}
The evaluation uses Mirai and BASHLITE variants on N-BaIoT.
No poisoning, backdoor, or evasion evaluation is performed.

\textbf{Privacy.}
B4's fingerprint constitutes distributional metadata leakage; B4 is not a privacy mechanism.
Concretely, the mean, variance, skewness, and $p_{95}$ reconstruction-error summaries may
reveal device-specific traffic regularity, tail heaviness, or unusual benign-error profiles;
secure aggregation or noisy aggregation would reduce this metadata exposure.
Mitigations (secure aggregation, differential privacy) are future work.

\textbf{Convergence and budget-capped runs.}
Budget-capped runs limit broad convergence interpretation, but within each seed the
same frozen AE is used across B1--B4, preserving threshold-calibration-only attribution.

\textbf{Architecture and $E{=}1$ scope.}
With $E{=}1$ and full participation, FedAvg approaches synchronous distributed training.
The results characterize threshold-scope personalization under a frequently
synchronized encoder; they should not be read as evidence that the same magnitude
holds under drift-heavy settings ($E>1$ is future work).

\textbf{Concept drift.}
This study uses fixed train/calibration/test splits and does not model online
recalibration; DATP assumes benign reconstruction-error distributions remain
sufficiently stable between calibration and deployment.

\textbf{Seeds and statistical power.}
Five seeds provide limited statistical power.
Seed-level descriptive summaries (all deltas positive, range reported) are appropriate
for this sample size.
All claims are scoped to the tested datasets, partitions, and protocol.
